\documentclass[../../main.tex]{subfiles}
\graphicspath{{../../image/}{../../figure/}}
\begin{document}

\section{弱解}\label{section:ch07-弱解}

鉴于上一节所展示的古典解的破裂，为在大范围求解非线性双曲守恒律组的初值问题，必须诉诸弱解。如第 IV 章所解释，弱解的可容许性问题必须加以处理。

在前几章中，我们主要考虑仅是有界可测函数的弱解。在该函数类中的存在性将在第 XVII 章中通过补偿紧性的泛函解析方法，对某些方程组加以确立。另一方面，存在三守恒律组的 Cauchy 问题在 $L^1$ 中不适定的情形。显然，一维空间中双曲守恒律组的选定函数类是 $BV$，它为设想弱解的最重要特征（即激波及其相互作用）提供了自然框架。

双曲方程组解的有限依赖域性质，结合我们的方程组\eqref{align:ch07-典范守恒律}在均匀伸缩
\begin{align}
x=\bar{x}+a y,\quad t=\bar{t}+a\tau,\quad a>0
\label{align:ch07-均匀伸缩}
\end{align}
下不变的事实，提示 $BV$ 弱解的可容许性可在局部通过考察激波与波扇来决定。这些问题将在接下来两章中充分讨论。

\end{document}
